\documentclass[../main.tex]{subfiles}
\ifSubfilesClassLoaded{\setcounter{chapter}{11}}{}
\begin{document}

\chapter{Berkas Teks dan Biner dengan API \texttt{FILE}}

\begin{subcpmk}
  \item Sub-CPMK 7.1: \texttt{fopen}/\texttt{fclose}, mode buka file, \texttt{FILE *}
  \item Sub-CPMK 7.2: Membaca/menulis file teks dan biner
  \item Sub-CPMK 7.3: Posisi file (\texttt{fseek}/\texttt{ftell}), penanganan error
\end{subcpmk}

\SesiRPS{12}{7.1--7.3}{$\sim$170 menit}{CPMK-7}

% ============================================================
\section{Konsep File I/O}
% ============================================================

File I/O (\textit{Input/Output}) memungkinkan program menyimpan data secara \textbf{persisten} --- data tetap ada bahkan setelah program berakhir. Dalam C, operasi file dilakukan melalui \textbf{stream}: aliran data antara program dan file di disk. Pustaka \texttt{<stdio.h>} menyediakan tipe \texttt{FILE} dan fungsi-fungsi untuk membuka, membaca, menulis, dan menutup file.

Setiap program C secara otomatis memiliki tiga stream yang sudah terbuka: \texttt{stdin} (input standar, biasanya keyboard), \texttt{stdout} (output standar, biasanya layar), dan \texttt{stderr} (output error). Fungsi \texttt{printf} sebenarnya menulis ke \texttt{stdout}, dan \texttt{scanf} membaca dari \texttt{stdin}. Dengan file I/O, kita menambahkan stream baru yang terhubung ke file di disk.

% ============================================================
\section{Membuka dan Menutup File}
% ============================================================

\begin{tabular}{|l|p{9cm}|}
\hline
\textbf{Mode} & \textbf{Keterangan} \\
\hline
\texttt{"r"} & Baca teks (\textit{read}). File harus sudah ada. \\
\hline
\texttt{"w"} & Tulis teks (\textit{write}). Membuat file baru atau \textbf{menimpa} file lama. \\
\hline
\texttt{"a"} & Tambah teks (\textit{append}). Menulis di akhir file. \\
\hline
\texttt{"r+"} & Baca dan tulis. File harus sudah ada. \\
\hline
\texttt{"w+"} & Baca dan tulis. Membuat baru / menimpa. \\
\hline
\texttt{"a+"} & Baca dan tambah di akhir. \\
\hline
\texttt{"rb"} & Baca biner. \\
\hline
\texttt{"wb"} & Tulis biner. \\
\hline
\end{tabular}

\begin{konsep}
\textbf{Pola standar File I/O:}

\begin{enumerate}[nosep]
  \item \texttt{fopen()} --- buka file, dapatkan \texttt{FILE *}
  \item Periksa apakah pointer \texttt{NULL} (file gagal dibuka)
  \item Baca atau tulis data
  \item \texttt{fclose()} --- \textbf{selalu} tutup file saat selesai
\end{enumerate}
\end{konsep}

% ============================================================
\section{Membaca File Teks}
% ============================================================

\begin{lstlisting}[language=C,caption=Membaca file baris per baris]
#include <stdio.h>

int main(void) {
    FILE *fp = fopen("data.txt", "r");
    
    if (fp == NULL) {
        perror("Gagal membuka file");  // cetak pesan error
        return 1;
    }
    
    char baris[256];
    int no = 1;
    
    while (fgets(baris, sizeof(baris), fp) != NULL) {
        printf("Baris %d: %s", no, baris);
        no++;
    }
    
    fclose(fp);  // WAJIB ditutup!
    return 0;
}
\end{lstlisting}

{\small
\begin{tabular}{|>{\raggedright\arraybackslash}p{3.0cm}|>{\raggedright\arraybackslash}p{3.2cm}|>{\raggedright\arraybackslash}p{6.5cm}|}
\hline
\textbf{Fungsi Baca} & \textbf{Unit} & \textbf{Catatan} \\
\hline
\texttt{fgetc(fp)} & 1 karakter & Mengembalikan \texttt{int} (\texttt{EOF} jika selesai) \\
\hline
\texttt{fgets(buf, n, fp)} & 1 baris (maks $n-1$ char) & Menyertakan \texttt{\textbackslash n}; aman \\
\hline
\texttt{fscanf(fp, fmt, ...)} & Terformat & Seperti \texttt{scanf} tapi dari file \\
\hline
\texttt{fread(buf, size, count, fp)} & Blok biner & Untuk file biner \\
\hline
\end{tabular}
}

% ============================================================
\section{Menulis ke File Teks}
% ============================================================

\begin{lstlisting}[language=C,caption=Menulis data ke file]
#include <stdio.h>

typedef struct {
    char nama[50];
    int skor;
} Siswa;

int main(void) {
    Siswa daftar[] = {
        {"Andi", 85}, {"Budi", 92}, {"Cici", 78}
    };
    int n = sizeof(daftar) / sizeof(daftar[0]);
    
    FILE *fp = fopen("hasil.txt", "w");
    if (!fp) { perror("Error"); return 1; }
    
    fprintf(fp, "%-20s %s\n", "Nama", "Skor");
    fprintf(fp, "-------------------------\n");
    for (int i = 0; i < n; i++) {
        fprintf(fp, "%-20s %d\n", daftar[i].nama, daftar[i].skor);
    }
    
    fclose(fp);
    printf("Data berhasil ditulis ke hasil.txt\n");
    return 0;
}
\end{lstlisting}

{\small
\begin{tabular}{|>{\raggedright\arraybackslash}p{3.0cm}|>{\raggedright\arraybackslash}p{3.2cm}|>{\raggedright\arraybackslash}p{6.5cm}|}
\hline
\textbf{Fungsi Tulis} & \textbf{Unit} & \textbf{Catatan} \\
\hline
\texttt{fputc(c, fp)} & 1 karakter & \\
\hline
\texttt{fputs(str, fp)} & 1 string & Tidak menambah newline otomatis \\
\hline
\texttt{fprintf(fp, fmt, ...)} & Terformat & Seperti \texttt{printf} tapi ke file \\
\hline
\texttt{fwrite(buf, size, count, fp)} & Blok biner & Untuk file biner \\
\hline
\end{tabular}
}

% ============================================================
\section{File Biner}
% ============================================================

File biner menyimpan data dalam format internal komputer (bukan teks yang dapat dibaca manusia). File biner lebih efisien untuk data terstruktur karena tidak perlu konversi format. Gunakan \texttt{fwrite} untuk menulis dan \texttt{fread} untuk membaca.

\begin{lstlisting}[language=C,caption=Menulis dan membaca struct ke/dari file biner]
#include <stdio.h>

typedef struct {
    char nama[30];
    double ipk;
} Mhs;

int main(void) {
    // TULIS ke file biner
    Mhs tulis[] = {{"Andi", 3.5}, {"Budi", 3.8}};
    FILE *fw = fopen("mhs.dat", "wb");
    if (!fw) { perror("Error"); return 1; }
    fwrite(tulis, sizeof(Mhs), 2, fw);
    fclose(fw);
    
    // BACA dari file biner
    Mhs baca[2];
    FILE *fr = fopen("mhs.dat", "rb");
    if (!fr) { perror("Error"); return 1; }
    size_t dibaca = fread(baca, sizeof(Mhs), 2, fr);
    fclose(fr);
    
    printf("Dibaca %zu record:\n", dibaca);
    for (size_t i = 0; i < dibaca; i++) {
        printf("  %s - %.2f\n", baca[i].nama, baca[i].ipk);
    }
    
    return 0;
}
\end{lstlisting}

% ============================================================
\section{Posisi File dan Penanganan Error}
% ============================================================

\begin{tabular}{|l|p{8cm}|}
\hline
\textbf{Fungsi} & \textbf{Keterangan} \\
\hline
\texttt{ftell(fp)} & Mengembalikan posisi saat ini (byte dari awal) \\
\hline
\texttt{fseek(fp, offset, origin)} & Pindah posisi. Origin: \texttt{SEEK\_SET} (awal), \texttt{SEEK\_CUR} (saat ini), \texttt{SEEK\_END} (akhir) \\
\hline
\texttt{rewind(fp)} & Kembali ke awal file \\
\hline
\texttt{feof(fp)} & Cek apakah sudah di akhir file \\
\hline
\texttt{ferror(fp)} & Cek apakah terjadi error pada stream \\
\hline
\texttt{perror(msg)} & Cetak pesan error berdasarkan \texttt{errno} \\
\hline
\end{tabular}

\begin{lstlisting}[language=C,caption=Menghitung ukuran file dengan \texttt{fseek}/\texttt{ftell}]
#include <stdio.h>

int main(void) {
    FILE *fp = fopen("data.txt", "r");
    if (!fp) { perror("Error"); return 1; }
    
    fseek(fp, 0, SEEK_END);     // pindah ke akhir
    long ukuran = ftell(fp);    // posisi = ukuran file
    rewind(fp);                 // kembali ke awal
    
    printf("Ukuran file: %ld byte\n", ukuran);
    fclose(fp);
    return 0;
}
\end{lstlisting}

% ============================================================
\section{Referensi sesi}
% ============================================================
\begin{itemize}[leftmargin=*]
  \item cppreference, I/O berkas. \url{https://en.cppreference.com/w/c/io}
  \item Programiz, \textit{C File Handling}. \url{https://www.programiz.com/c-programming/c-file-input-output}
  \item Petani Kode, file. \url{https://petanikode.com/tutorial/c}
\end{itemize}

% ============================================================
\begin{aktivitas}
  \item Buat program yang membaca file teks dan menghitung jumlah baris, kata, dan karakter.
  \item Buat program yang menulis array of struct ke file teks dengan format tabel.
  \item Buat program yang menyimpan dan memuat data struct ke/dari file biner.
\end{aktivitas}

\begin{latihan}
  \item Apa yang terjadi jika \texttt{fclose} tidak dipanggil?
  \item Jelaskan perbedaan mode \texttt{"w"} dan \texttt{"a"}.
  \item Kapan file biner lebih tepat daripada file teks?
  \item Apa kegunaan \texttt{perror()} dan bagaimana ia bekerja?
\end{latihan}

\begin{asesmen}
\textbf{Tugas M10:} program yang membaca file teks, menghitung jumlah baris, dan menulis ringkasan ke file output baru.
\end{asesmen}

\begin{checklist}
  \item Saya selalu memeriksa nilai balik \texttt{fopen} (NULL = gagal)
  \item Saya selalu menutup file dengan \texttt{fclose}
  \item Saya dapat membaca dan menulis file teks
  \item Saya memahami perbedaan file teks dan file biner
  \item Saya dapat menggunakan \texttt{fseek}/\texttt{ftell} untuk navigasi file
\end{checklist}

\begin{rangkuman}
Bab ini membahas file I/O dalam C: membuka file (\texttt{fopen} dengan berbagai mode), membaca (\texttt{fgetc}, \texttt{fgets}, \texttt{fscanf}, \texttt{fread}), menulis (\texttt{fputc}, \texttt{fputs}, \texttt{fprintf}, \texttt{fwrite}), menutup file (\texttt{fclose}), navigasi file (\texttt{fseek}, \texttt{ftell}), dan penanganan error (\texttt{perror}).

\textbf{Poin kunci:} selalu cek \texttt{fopen} NULL; selalu \texttt{fclose}; mode \texttt{"w"} menimpa file; \texttt{fgets} lebih aman dari \texttt{fscanf} untuk baris; file biner untuk data terstruktur.

\textbf{Kata kunci:} \texttt{FILE *}, \texttt{fopen}, \texttt{fclose}, \texttt{fgets}, \texttt{fprintf}, \texttt{fread}, \texttt{fwrite}, \texttt{fseek}, mode buka
\end{rangkuman}

\end{document}
